\caption{\textbf{Field or coupling: activity co-moves through the scheduler, talk spreads through one read-out call at a time.}
(a)~Simulation on a synthetic call schedule (ticks: model calls, $\sim$13\,s apart). A field (orange line) reaches every agent at its next call, so all respond within about one call. A coupling passes from reader to reader, and each hop waits for the recipient's next call.
(b)~G38 (regime III), measured lags. Orange: each agent's first call of the day relative to that day's median (pooled IQR 12\,s; per-day median IQR 9\,s; 82\% of agents start before any peer message reaches them). Blue: from a peer message to the recipient's receiving call (median 17\,s; regime-III unit median 24\,s). The tail is pauses.
(c)~Excess chance that a recipient's call is a talk call, by hop after a peer message (IVW pooled over 41 regime-I and 27 regime-III units; bands are 95\% CIs). It is zero at the call in flight and steps up at the read-out call. The gray band is the pooled shifted-time placebo. $J_1>0$ in 45/71 units and $<0$ in none. The regime-I rise beyond hop 2 is not validated on synthetic data.}
\label{fig:H50}
